\chapter{Podstawy algebry abstrakcyjnej: grupy, pierścienie i moduły}
\label{ch:algebra-abstrakcyjna-v5}

W topologii algebraicznej każdej przestrzeni przypisujemy grupy, a odwzorowaniom
ciągłym --- homomorfizmy. Grupa podstawowa może być nieprzemienna, grupy
homologii są abelowe, a iloczyn tensorowy i funktory $\operatorname{Tor}$,
$\operatorname{Ext}$ opisują oddziaływanie współczynników. Dlatego zbieramy
potrzebne konstrukcje \emph{przed} rozdziałem o homologii algebraicznej.
Żadne twierdzenie niniejszego rozdziału nie odsyła do osobnego rozdziału
książki geometrycznej o algebrze tensorowej. Pracujemy nad pierścieniami
z jedynką, homomorfizmy pierścieni zachowują jedynkę, a moduły są lewostronne,
jeżeli nie wskazano inaczej.

\section{Grupy, homomorfizmy, podgrupy normalne}
\begin{definicja}
\emph{Grupa} to zbiór $G$ z działaniem $(g,h)\mapsto gh$ łącznym,
elementem neutralnym $1_G$ i odwrotnościami $g^{-1}$ wszystkich elementów.
Jest \emph{abelowa}, gdy $gh=hg$ dla wszystkich $g,h\in G$.
\emph{Homomorfizm} $\varphi:G\to H$ spełnia
$\varphi(gh)=\varphi(g)\varphi(h)$. Podgrupa $N\leq G$ jest
\emph{normalna}, $N\trianglelefteq G$, gdy $gNg^{-1}=N$ dla wszystkich
$g\in G$. Jądro homomorfizmu to $\ker\varphi=\varphi^{-1}(\{1_H\})$.
\end{definicja}
\begin{przyklad}
Dla $m\geq1$ grupy $(\Z,+)$, $(\Z/m\Z,+)$,
grupa permutacji $S_m$ i grupa macierzy $\mathrm{GL}_m(\R)$
są grupami. Pierwsze dwie są abelowe;
$S_m$ jest nieabelowa dla $m\geq3$, a
$\mathrm{GL}_m(\R)$ dla $m\geq2$. Homomorfizm $\Z\to \Z/m\Z$ wysyłający $k$ na jego
resztę ma jądro $m\Z$. Wyznacznik jest homomorfizmem
$\det:\mathrm{GL}_m(\R)\to\R^\times$ z jądrem
$\mathrm{SL}_m(\R)$.
\end{przyklad}
\begin{stwierdzenie}
Dla każdego homomorfizmu grup $\varphi:G\to H$ jego jądro jest
normalną podgrupą $G$, a obraz jest podgrupą $H$.
\end{stwierdzenie}
\begin{proof}
Jeżeli $n\in\ker\varphi$ i $g\in G$, to
$\varphi(gng^{-1})=\varphi(g)1_H\varphi(g)^{-1}=1_H$.
Pozostałe aksjomaty wynikają z zachowywania iloczynu, jedynki
i odwrotności przez $\varphi$.
\end{proof}

\begin{definicja}
Dla $N\trianglelefteq G$ na zbiorze klas lewostronnych $G/N$ definiujemy
$(gN)(hN)=ghN$. Tak otrzymana grupa nazywa się \emph{ilorazową}.
\end{definicja}
\begin{lemat}
Iloczyn na $G/N$ jest dobrze określony wtedy i tylko wtedy, gdy $N$
jest normalna. Naturalny rzut $q:G\to G/N$ jest homomorfizmem.
\end{lemat}
\begin{proof}
Przy normalności $g'=gn$, $h'=hm$ implikuje
$g'h'=gh(h^{-1}nh)m$, zatem $g'h'N=ghN$. W drugą stronę,
niezależność iloczynu od reprezentantów daje
$(gN)(nN)(g^{-1}N)=(gN)(N)(g^{-1}N)=N$,
więc $gng^{-1}\in N$. Włączenie dla $g^{-1}$ daje równość.
\end{proof}
\begin{twierdzenie}[Pierwsze twierdzenie o izomorfizmie]
\label{thm:izo-grupy-v5}
Jeżeli $\varphi:G\to H$ jest homomorfizmem grup, to
$G/\ker\varphi\cong\operatorname{im}\varphi$, przy czym izomorfizm
wysyła $g\ker\varphi$ na $\varphi(g)$.
\end{twierdzenie}
\begin{proof}
Jeśli $g'=gk$ dla $k\in\ker\varphi$, to
$\varphi(g')=\varphi(g)\varphi(k)=\varphi(g)$,
więc wzór nie zależy od przedstawiciela warstwy.
Jeżeli $\varphi(g)=\varphi(h)$, to
$h^{-1}g\in\ker\varphi$, a zatem
$g\ker\varphi=h\ker\varphi$, co dowodzi różnowartościowości.
Surjektywność na obraz wynika z jego definicji, a zgodność
z mnożeniem z faktu, że $\varphi$ jest homomorfizmem.
\end{proof}
\begin{uwaga}
W topologii, dla łukowo spójnej przestrzeni $X$ i punktu
$x_0\in X$, otrzymamy w szczególności naturalną surjekcję
$\pi_1(X,x_0)\twoheadrightarrow H_1(X;\Z)$, której jądrem jest
podgrupa komutatorowa grupy podstawowej. Jeżeli $X$ nie jest łukowo
spójna, analogiczne stwierdzenie dotyczy grupy $H_1(C;\Z)$
składowej łukowej $C$ zawierającej $x_0$, a nie automatycznie
całej grupy $H_1(X;\Z)$.
\end{uwaga}

\section{Działania grup, orbity, stabilizatory i iloczyny półproste}
\begin{definicja}
\emph{Lewym działaniem} grupy $G$ na zbiorze $X$ nazywamy funkcję
$(g,x)\mapsto gx$ spełniającą $1_Gx=x$ i $(gh)x=g(hx)$.
\emph{Orbita} $x$ to $Gx=\{gx:g\in G\}$, a \emph{stabilizator}
$G_x=\{g\in G:gx=x\}$.
\end{definicja}
\begin{twierdzenie}[Orbita i stabilizator]
Dla działania $G\curvearrowright X$ i każdego $x\in X$ istnieje
bijekcja zbiorów $G/G_x\to Gx$, $gG_x\mapsto gx$.
Tutaj $G/G_x$ oznacza zbiór warstw lewostronnych; stabilizator
$G_x$ nie musi być podgrupą normalną.
Jeżeli $G$ jest skończona, to $|Gx|=[G:G_x]$.
\end{twierdzenie}
\begin{proof}
Równość $gx=hx$ jest równoważna $h^{-1}g\in G_x$,
czyli $gG_x=hG_x$. Wynika stąd, że przyporządkowanie jest dobrze określone
i różnowartościowe. Każdy punkt orbity ma postać $gx$,
zatem jest ono również surjektywne.
\end{proof}
\begin{przyklad}
Działanie $\Z$ na $\R$ przez $n\cdot t=t+n$ ma iloraz
$\R/\Z\cong S^1$; stabilizatory są trywialne.
Dla $m\geq2$ działanie $S_m$ na $\{1,\ldots,m\}$
jest przechodnie, a stabilizator $1$ jest izomorficzny
z $S_{m-1}$.
\end{przyklad}

\begin{definicja}[Iloczyn prosty i półprosty]
Dla grup $N,H$ iloczyn prosty $N\times H$ ma mnożenie
$(n,h)(n',h')=(nn',hh')$. Jeżeli
$\alpha:H\to\operatorname{Aut}(N)$ jest homomorfizmem,
to \emph{iloczyn półprosty} $N\rtimes_\alpha H$ jest zbiorem
$N\times H$ z mnożeniem
$(n,h)(n',h')=(n\alpha_h(n'),hh')$.
\end{definicja}
\begin{stwierdzenie}
Powyższe mnożenie jest grupowe; elementem neutralnym jest
$(1_N,1_H)$, a odwrotność $(n,h)$ wynosi
$(\alpha_{h^{-1}}(n^{-1}),h^{-1})$. Jeśli $\alpha$ jest
trywialne, otrzymujemy iloczyn prosty.
\end{stwierdzenie}
\begin{proof}
Oznaczmy trzy elementy przez $(n,h)$, $(n',h')$, $(n'',h'')$.
Z pierwszego sposobu ustawienia nawiasów otrzymujemy
pierwszą współrzędną
$n\alpha_h(n')\alpha_{hh'}(n'')$.
Z drugiego otrzymujemy
$n\alpha_h(n'\alpha_{h'}(n''))$, co daje ten sam wynik,
ponieważ $\alpha_h$ jest automorfizmem i
$\alpha_h\alpha_{h'}=\alpha_{hh'}$.
Druga współrzędna jest w obu przypadkach $hh'h''$.
Ponadto $\alpha_{1_H}=\id_N$, więc $(1_N,1_H)$
jest elementem neutralnym.
Dla $b=(\alpha_{h^{-1}}(n^{-1}),h^{-1})$ zachodzi
$(n,h)b=(n\alpha_h\alpha_{h^{-1}}(n^{-1}),1_H)
=(1_N,1_H)$.
W odwrotnej kolejności
$b(n,h)=(\alpha_{h^{-1}}(n^{-1}n),1_H)=(1_N,1_H)$.
Dostajemy grupę. Gdy $\alpha_h=\id_N$ dla wszystkich $h$,
wzór sprowadza się do iloczynu prostego.
\end{proof}
\begin{przyklad}[Grupa butelki Kleina]
Niech $H=\Z$, $N=\Z$ i $\alpha_k(n)=(-1)^k n$. Wtedy
$N\rtimes_\alpha H$ jest grupą o prezentacji
$\langle a,b\mid aba^{-1}=b^{-1}\rangle$.
Istotnie, dla $a=(0,1)$ i $b=(1,0)$ mamy
$aba^{-1}=b^{-1}$; każdy $(n,k)$ ma postać $b^na^k$.
Z samej relacji $ab=b^{-1}a$ wynika, że każde słowo
w prezentowanej grupie można sprowadzić do takiej postaci.
Ponieważ w iloczynie półprostym wszystkie pary $(n,k)$
są różne, żadne dodatkowe relacje nie zachodzą.
To pojawi się przy obliczaniu grupy podstawowej butelki Kleina
z monodromii fibracji okręgowej.
\end{przyklad}

\begin{przyklad}[Symetrie kwadratu i iloczyn półprosty]
\label{prz:dihedral-top-algebra}
Niech $r$ będzie obrotem kwadratu o kąt $\pi/2$,
a $s$ odbiciem względem osi pionowej.
Dla grupy symetrii $D_4$ mamy
$r^4=s^2=1$ oraz $srs=r^{-1}$.
Każda symetria daje się jednoznacznie przedstawić
w postaci $r^k$ lub $r^ks$, gdzie $k=0,1,2,3$:
obrazy wyróżnionego wierzchołka i jego sąsiada
wyznaczają symetrię, a obrotów i odbić jest po cztery.
Dlatego $D_4\cong(\Z/4\Z)\rtimes(\Z/2\Z)$,
gdzie niezerowy element drugiego czynnika działa
na pierwszym przez $k\mapsto-k$.
Działanie $D_4$ na czterech wierzchołkach kwadratu
jest przechodnie, a stabilizator jednego wierzchołka
ma dwa elementy (identyczność i odbicie przez przekątną).
Z twierdzenia o orbicie i stabilizatorze otrzymujemy
$4=8/2$.
\end{przyklad}

\begin{figure}[H]
\centering
\begin{tikzpicture}[>=Latex,scale=1.0,font=\small]
  \coordinate (A) at (-1.15,1.05);
  \coordinate (B) at (1.15,1.05);
  \coordinate (C) at (1.15,-1.05);
  \coordinate (D) at (-1.15,-1.05);
  \draw[very thick] (A)--(B)--(C)--(D)--cycle;
  \foreach \P/\name/\dir in {A/1/above left,B/2/above right,
                            C/3/below right,D/4/below left}
    {\fill (\P) circle(2.2pt); \node[\dir] at (\P) {$\name$};}
  \draw[->,manifoldblue,very thick] (-.536,.45)
    arc[start angle=140,end angle=40,radius=.70];
  \node[manifoldblue] at (.72,.05) {$r$};
  \draw[manifoldred,dashed,thick] (0,-1.38)--(0,1.40);
  \node[manifoldred,anchor=west] at (.08,1.42) {$s$};
\end{tikzpicture}
\caption{Dwa generatory grupy symetrii kwadratu:
obrót $r$ i odbicie $s$. Wybór orientacji obrotu
nie zmienia relacji $srs=r^{-1}$.}
\label{fig:dihedral-top-algebra}
\end{figure}

\section{Grupy wolne i prezentacje przez generatory i relacje}
\begin{definicja}
Niech $S$ będzie zbiorem. \emph{Grupa wolna} $F(S)$ składa się ze
zredukowanych skończonych słów w symbolach $s,s^{-1}$, $s\in S$,
z mnożeniem przez konkatenację i redukcję par $ss^{-1}$ i $s^{-1}s$.
\emph{Prezentacja} $\langle S\mid R\rangle$ oznacza iloraz $F(S)$
przez najmniejszą normalną podgrupę zawierającą słowa relacji $R$.
Dla $S=\varnothing$ otrzymujemy grupę trywialną.
\end{definicja}

\begin{uwaga}[Jednoznaczność słowa zredukowanego]
Aby konstrukcja $F(S)$ była ścisła, trzeba sprawdzić,
że wynik redukcji nie zależy od kolejności usuwania par.
Każde skreślenie zmniejsza długość słowa o dwa, więc
proces kończy się. Dwa możliwe pierwsze skreślenia
albo dotyczą rozłącznych par i mogą być wykonane
w dowolnej kolejności, albo zachodzą na siebie,
jak w $ss^{-1}s$ lub $s^{-1}ss^{-1}$; wówczas
oba dają to samo krótsze słowo. W pierwszym przypadku
oba skreślenia prowadzą po następnym kroku do tego samego
słowa, a w drugim już po jednym kroku. Indukcja po długości
słowa dowodzi zatem, że każda kolejność skreśleń kończy się
tym samym słowem zredukowanym.
Iloczyn słów definiujemy jako zredukowaną konkatenację.
Jeżeli $\rho(w)$ oznacza jednoznaczną redukcję słowa $w$, to
$\rho(\rho(uv)w)=\rho(uvw)=\rho(u\rho(vw))$,
co dowodzi łączności. Elementem neutralnym jest słowo puste,
a odwrotność otrzymujemy, odwracając kolejność liter
i zastępując każdą z nich jej odwrotnością.
\end{uwaga}
\begin{twierdzenie}[Własność uniwersalna grupy wolnej]
Dla dowolnej grupy $H$ i funkcji $u:S\to H$ istnieje dokładnie
jeden homomorfizm $\overline u:F(S)\to H$ przedłużający $u$.
Homomorfizm ten przechodzi na $\langle S\mid R\rangle$ wtedy
i tylko wtedy, gdy wszystkie relacje z $R$ mają obraz $1_H$.
\end{twierdzenie}
\begin{proof}
Definiujemy obraz słowa, mnożąc obrazy kolejnych liter,
a dla $s^{-1}$ biorąc $u(s)^{-1}$. Usunięcie sąsiednich
liter odwrotnych nie zmienia wartości. Każdy homomorfizm
musi działać na słowach w ten sposób, co daje jednoznaczność.
Mapa przechodzi na iloraz dokładnie wtedy, gdy
jej jądro zawiera najmniejszą normalną podgrupę
generowaną przez $R$. Ponieważ jądro jest normalne,
jest to równoważne warunkowi, że każde słowo relacji
jest wysyłane na element neutralny $H$.
\end{proof}
\begin{przyklad}
Grupa podstawowa bukietu $r$ okręgów jest wolna na $r$
generatorach. Dołączenie $2$-komórki wzdłuż słowa
$aba^{-1}b^{-1}$ wymusza przemienność generatorów i daje
$\langle a,b\mid aba^{-1}b^{-1}\rangle\cong\Z^2$.
Ścisły związek z dołączaniem komórek dowiedziemy z twierdzenia
van Kampena.
\end{przyklad}

\section{Pierścienie, ideały i konstrukcja ilorazu}
\begin{definicja}
\emph{Pierścień z jedynką} $(R,+,\cdot,0,1)$ ma abelową grupę
addytywną, łączne mnożenie z jedynką oraz rozdzielność mnożenia
względem dodawania. Jest \emph{przemienny}, gdy $ab=ba$.
Podzbiór $I\subseteq R$ jest \emph{ideałem dwustronnym}, gdy
jest podgrupą addytywną i $rI,Ir\subseteq I$ dla wszystkich $r\in R$.
\end{definicja}
\begin{przyklad}
$\Z$, $\Q[x]$, pierścień macierzy $M_n(\R)$ i pierścień grupowy
$\Z[G]$ są pierścieniami z jedynką. Pierścień $M_n(\R)$
przy $n\ge2$ jest nieprzemienny. W $\Z$ każdy ideał ma postać
$m\Z$; w $\Q[x]$ ideały mają postać $(f)$.
\end{przyklad}
\begin{twierdzenie}[Pierścień ilorazowy]
Jeżeli $I\triangleleft R$ jest ideałem dwustronnym, to na $R/I$
wzory $(r+I)+(s+I)=(r+s)+I$ i $(r+I)(s+I)=rs+I$
definiują pierścień. Jeżeli $\varphi:R\to S$ jest homomorfizmem
pierścieni, to $R/\ker\varphi\cong\operatorname{im}\varphi$.
\end{twierdzenie}
\begin{proof}
Dla $r'=r+i$, $s'=s+j$ mamy
$r's'-rs=rj+is+ij\in I$, stąd mnożenie nie zależy
od reprezentantów. Aksjomaty dziedziczymy z $R$.
Izomorfizm ilorazu z obrazem jest zadany przez $r+\ker\varphi
\mapsto\varphi(r)$; sprawdzenie jego poprawności i bijektywności
przebiega tak jak dla grup.
\end{proof}
\begin{definicja}
\emph{Dziedziną całkowitości} nazywamy niezerowy pierścień
przemienny z jedynką bez dzielników zera. \emph{Ciałem} jest
niezerowy pierścień przemienny z jedynką, w którym każdy
element różny od zera jest odwracalny. \emph{Dziedziną ideałów głównych} (PID) jest
dziedzina, w której każdy ideał jest generowany przez jeden element.
\end{definicja}
\begin{przyklad}
$\Z$ jest PID na mocy dzielenia z resztą, podobnie $K[x]$
dla dowolnego ciała $K$. Pierścień $\Z[x]$ nie jest PID:
ideał $(2,x)$ nie jest główny. Gdyby $(2,x)=(f)$,
to $f$ musiałby dzielić zarówno $2$, jak i $x$. Ponieważ
$f$ dzieli stały wielomian $2$, sam musi być stały i równy
$\pm1$ lub $\pm2$. Stała $\pm2$ nie dzieli $x$, więc
$f$ byłby jednostką i $(2,x)=\Z[x]$. Tymczasem
$1\notin(2,x)$: po podstawieniu $x=0$ i redukcji modulo $2$
każdy element tego ideału daje zero. Pierścień $\Z/6\Z$ ma dzielniki zera,
ponieważ $[2][3]=[0]$, więc nie jest dziedziną.
\end{przyklad}

\section{Moduły, podmoduły, ilorazy i ranga}
\begin{definicja}
Niech $R$ będzie pierścieniem z jedynką. \emph{Lewy $R$-moduł}
$M$ jest grupą abelową z działaniem $R\times M\to M$ takim,
że $(rs)m=r(sm)$, $1m=m$ oraz zachodzą obie rozdzielności.
Podmoduł $N\subseteq M$ jest podgrupą addytywną zamkniętą
na mnożenie przez skalar. Homomorfizm $R$-modułów zachowuje
sumy i mnożenie przez elementy $R$.
\end{definicja}
\begin{przyklad}
Moduły nad $\Z$ to dokładnie grupy abelowe, a moduły nad
ciałem to przestrzenie wektorowe. $2\Z\subset\Z$ jest
podmodułem, którego ilorazem jest $\Z/2\Z$.
Przy pierścieniu $R=\Z/4\Z$ podmoduł $2R=\{[0],[2]\}$
nie jest wolnym $R$-modułem. Gdyby miał niepustą bazę,
odwzorowanie $R\to2R$, $r\mapsto re$, wyznaczone przez
dowolny element bazowy $e$, byłoby różnowartościowe.
Jest to niemożliwe, bo $R$ ma cztery elementy,
a $2R$ tylko dwa. Baza pusta również odpada,
ponieważ $2R\ne0$.
\end{przyklad}
\begin{twierdzenie}[Pierwsze twierdzenie o izomorfizmie modułów]
Jeżeli $f:M\to N$ jest homomorfizmem $R$-modułów, to
$M/\ker f\cong\operatorname{im}f$ jako $R$-moduły.
\end{twierdzenie}
\begin{proof}
Jeżeli $m+\ker f=m'+\ker f$, to
$m-m'\in\ker f$, czyli $f(m)=f(m')$; wzór
$m+\ker f\mapsto f(m)$ jest więc dobrze określony.
Odwrotnie, równość $f(m)=f(m')$ oznacza
$m-m'\in\ker f$, co daje różnowartościowość.
Surjektywność na $\operatorname{im}f$ wynika z definicji obrazu,
a $f(m+m')=f(m)+f(m')$ oraz $f(rm)=rf(m)$
dowodzą liniowości indukowanego odwzorowania.
\end{proof}

\begin{definicja}
Rodzina $(e_i)_{i\in I}$ jest \emph{bazą modułu} $M$,
jeżeli każdy element ma jednoznaczne przedstawienie jako
skończona kombinacja $\sum_{i\in I}r_ie_i$.
Moduł z bazą nazywamy \emph{wolnym}. Zapis $R^{(I)}$
oznacza sumę prostą $|I|$ kopii $R$, a nie iloczyn wszystkich
kopii. \emph{Ciąg dokładny} $M'\xrightarrow{u}M\xrightarrow{v}M''$
spełnia $\operatorname{im}u=\ker v$; \emph{krótki ciąg dokładny}
ma postać $0\to M'\xrightarrow{u}M\xrightarrow{v}M''\to0$.
\end{definicja}
\begin{twierdzenie}[Rozszczepienie krótkiego ciągu]
Dla krótkiego ciągu dokładnego $R$-modułów
$0\to A\xrightarrow{i}B\xrightarrow{p}C\to0$
następujące warunki są równoważne:
(i) istnieje $R$-liniowe $s:C\to B$ z $ps=\id_C$;
(ii) istnieje $R$-liniowe $r:B\to A$ z $ri=\id_A$;
(iii) istnieje izomorfizm $R$-modułów $\theta:B\to A\oplus C$
taki, że $\theta(i(a))=(a,0)$ dla każdego $a\in A$ oraz
$\operatorname{pr}_C\circ\theta=p$.
\end{twierdzenie}
\begin{proof}
Z (i) wynika, że $b=(b-sp(b))+sp(b)$, gdzie pierwszy
składnik należy do $\ker p=\operatorname{im}i$.
Suma jest prosta: jeśli $i(a)=s(c)$, to zastosowanie $p$
daje $c=0$, następnie $a=0$. Izomorfizm z (iii) jest więc
zadany przez $\theta(b)=(i^{-1}(b-sp(b)),p(b))$,
gdzie $i^{-1}$ oznacza odwrotność $i:A\to\operatorname{im}i$.
Z (iii) otrzymujemy naturalne włączenie $C\to A\oplus C$
i rzut na $A$, co daje (i) i (ii). Alternatywnie z (ii) mamy
$B=i(A)\oplus\ker r$, ponieważ każdy $b\in B$
rozkłada się jednoznacznie jako
$b=i(r(b))+(b-i(r(b)))$.
Ograniczenie $p|_{\ker r}$ jest różnowartościowe:
jeśli $b\in\ker r\cap\ker p$, to $b=i(a)$,
więc $a=r(i(a))=0$. Jest też surjektywne:
dla $c\in C$ wybieramy $b$ z $p(b)=c$,
a wtedy $b-i(r(b))\in\ker r$ ma obraz $c$.
Otrzymujemy więc izomorfizm $\ker r\cong C$.
W szczególności odwzorowanie
$\theta:B\to A\oplus C$, $\theta(b)=(r(b),p(b))$,
jest izomorfizmem: różnowartościowość wynika z
$\ker r\cap\ker p=\{0\}$, a surjektywność z rozkładu
$B=i(A)\oplus\ker r$ i bijektywności $p|_{\ker r}$.
Ponadto $\theta(i(a))=(a,0)$ i
$\operatorname{pr}_C\circ\theta=p$, zgodnie z (iii).
\end{proof}
\begin{przyklad}[Ciąg nierozszczepialny]
$0\to\Z\xrightarrow{\cdot 2}\Z\to\Z/2\Z\to0$
jest dokładny, ale nie rozszczepia się. Gdyby $s:\Z/2\Z\to\Z$
był przekrojem addytywnym, obraz klasy $[1]$ musiałby być
jednocześnie niezerowy (bo $ps([1])=[1]$) i mieć rząd $2$,
co jest niemożliwe w $\Z$.
\end{przyklad}


\section{Iloczyn tensorowy modułów: konstrukcja bez zewnętrznych odsyłaczy}
\label{sec:tensor-algebra-v5}
Dotychczas rozpatrywaliśmy moduły $M$ i $N$ osobno. Iloczyn
tensorowy reprezentuje odwzorowania \emph{dwuliniowe}:
zamiast rozpatrywać funkcję dwu zmiennych, można rozpatrywać
jedno odwzorowanie liniowe z nowego modułu.

\begin{definicja}[Odwzorowanie dwuliniowe]
Niech $R$ będzie pierścieniem przemiennym z jedynką,
a $M,N,P$ modułami nad $R$. Odwzorowanie
$b:M\times N\to P$ jest $R$-\emph{dwuliniowe}, jeśli jest
addytywne w każdej zmiennej oraz
$b(rm,n)=r b(m,n)=b(m,rn)$ dla wszystkich $r\in R$.
\end{definicja}
\begin{twierdzenie}[Istnienie i uniwersalność tensoru]
\label{thm:tensor-uni-v5}
Niech $R$ będzie przemiennym pierścieniem z jedynką,
a $M,N$ modułami nad $R$. Istnieje $R$-moduł $M\otimes_RN$
oraz odwzorowanie dwuliniowe
$\otimes:M\times N\to M\otimes_RN$ takie, że
dla każdego $R$-modułu $P$ i każdego odwzorowania
dwuliniowego $b:M\times N\to P$
istnieje dokładnie jedno faktoryzowanie
$b=\overline b\circ\otimes$ przez $R$-liniowe
$\overline b:M\otimes_RN\to P$.
Para ta jest jednoznaczna z dokładnością do jednoznacznego
izomorfizmu zachowującego elementy $m\otimes n$.
\end{twierdzenie}
\begin{proof}
Niech $F$ będzie wolnym $R$-modułem o bazie symboli $[m,n]$
indeksowanych przez pary $(m,n)\in M\times N$.
Niech $L$ będzie podmodułem generowanym przez elementy
\[
[m+m',n]-[m,n]-[m',n],\quad
[m,n+n']-[m,n]-[m,n'],\quad
[rm,n]-r[m,n],\quad[m,rn]-r[m,n].
\]
Definiujemy $M\otimes_RN=F/L$ i $m\otimes n=[m,n]+L$.
Relacje gwarantują dwuliniowość. Dla $b$ jak w tezie
wolność $F$ daje liniową mapę $B:F\to P$, $B([m,n])=b(m,n)$.
Dwuliniowość oznacza $L\subseteq\ker B$, więc $B$ przechodzi
na jedyną mapę $\overline b:F/L\to P$. Jedyność bierze się
stąd, że elementy $m\otimes n$ generują $F/L$.
Jeżeli $(T,t)$ ma tę samą własność, to uniwersalność daje
mapy $T\to M\otimes_RN$ i odwrotną. Sprawdźmy, że obrazy
$t(m,n)$ generują $T$: jeżeli $T'$ jest podmodułem przez nie
generowanym, to $t$ przyjmuje wartości w $T'$, więc
uniwersalność daje mapę $u:T\to T'$ przedłużającą $t$.
Jej złożenie z inkluzją $T'\hookrightarrow T$ jest
identycznością na $t$; jednoznaczność wymusza więc
$T'=T$. Obie kompozycje powyższych map są zatem
identycznościami na zbiorze generatorów i są identycznościami.
\end{proof}
\begin{przyklad}
Dla grup abelowych i $m,n\ge1$ zachodzi
$(\Z/m\Z)\otimes_\Z(\Z/n\Z)\cong\Z/d\Z$,
gdzie $d=\gcd(m,n)$. Wykażmy to bez pomijania
relacji. Każdy tensor prosty jest wielokrotnością
$t=[1]\otimes[1]$, bo
$[a]\otimes[b]=ab\,t$; zatem cały iloczyn
tensorowy jest generowany przez $t$.
Z $m[1]=0$ i $n[1]=0$ wynika $mt=nt=0$.
Tożsamość Bézouta daje $d=um+vn$ dla pewnych
$u,v\in\Z$, więc również $dt=0$.
Stąd mamy surjekcję
$\Z/d\Z\to(\Z/m\Z)\otimes_\Z(\Z/n\Z)$,
$[k]\mapsto kt$.

Odwzorowanie
$b:(\Z/m\Z)\times(\Z/n\Z)\to\Z/d\Z$,
$b([a],[b])=[ab]$, jest dobrze określone:
zmiana $a$ o wielokrotność $m$ lub $b$ o
wielokrotność $n$ nie zmienia iloczynu modulo $d$,
ponieważ $d$ dzieli $m$ i $n$.
Jest ono dwuliniowe; przez własność uniwersalną
indukuje homomorfizm iloczynu tensorowego na
$\Z/d\Z$ wysyłający $t$ na $[1]$.
Stanowi on odwrotność poprzedniej surjekcji.
W szczególności $\Z/2\Z\otimes_\Z\Z/3\Z=0$,
choć oba czynniki są niezerowe.
\end{przyklad}
\begin{stwierdzenie}[Rachunek tensorów]
Dla przemiennego pierścienia $R$ z jedynką i
$R$-modułów $M,N,P$ istnieją naturalne izomorfizmy
\[
R\otimes_RM\cong M,\qquad
(M\oplus N)\otimes_RP\cong(M\otimes_RP)\oplus(N\otimes_RP),
\qquad (M\otimes_RN)\otimes_RP\cong M\otimes_R(N\otimes_RP).
\]
\end{stwierdzenie}
\begin{proof}
Dla pierwszego izomorfizmu wzór $r\otimes m\mapsto rm$
jest liniowy z odwrotnością $m\mapsto1\otimes m$.
W drugim przypadku dwuliniowy wzór
$(m,n)\otimes p\mapsto(m\otimes p,n\otimes p)$
wyznacza homomorfizm do sumy prostej.
Odwrotność definiujemy na obu składnikach sumy:
$m\otimes p\mapsto(m,0)\otimes p$ oraz
$n\otimes p\mapsto(0,n)\otimes p$.
Własność uniwersalna zapewnia poprawność obu odwzorowań;
ich złożenia są identycznościami na tensorach prostych.
Dla trzeciego izomorfizmu dwuliniowość względem
$m,n$ pozwala, dla ustalonego $p$, przejść od
$(m,n)\mapsto m\otimes(n\otimes p)$
do mapy liniowej określonej na $M\otimes_RN$.
Zależność od $p$ jest liniowa, co sprawdzamy na
tensorach prostych, a następnie na ich sumach.
Ponowne zastosowanie uniwersalności daje
$(m\otimes n)\otimes p\mapsto m\otimes(n\otimes p)$.
Analogicznie konstruujemy mapę odwrotną; oba złożenia
są identycznościami na tensorach prostych,
które generują odpowiednie moduły.
\end{proof}
\begin{twierdzenie}[Prawostronna dokładność iloczynu tensorowego]
Jeżeli $M'\xrightarrow{u}M\xrightarrow{v}M''\to0$ jest
dokładnym ciągiem $R$-modułów nad przemiennym $R$,
to dla każdego $R$-modułu $N$ ciąg
$M'\otimes_RN\xrightarrow{u\otimes1}M\otimes_RN
\xrightarrow{v\otimes1}M''\otimes_RN\to0$ jest dokładny.
\end{twierdzenie}
\begin{proof}
Surjektywność $v\otimes1$ wynika z tego, że
każdy tensor prosty $m''\otimes n$ ma podniesienie
$m\otimes n$, ponieważ $v$ jest surjektywne.
Niech $Q=(M\otimes_RN)/\operatorname{im}(u\otimes1)$.
Ponieważ $v\circ u=0$, odwzorowanie $v\otimes1$
indukuje surjekcję $\beta:Q\to M''\otimes_RN$.
Z równości $\ker v=\operatorname{im}u$ wynika, że
wzór $(v(m),n)\mapsto[m\otimes n]\in Q$ nie zależy
od wyboru $m$: dwie możliwości różnią się o $u(m')$,
a $u(m')\otimes n$ ma zerową klasę w $Q$.
Wzór jest dwuliniowy, gdyż działania w $Q$ pochodzą
z $M\otimes_RN$. Uniwersalność daje więc
$\alpha:M''\otimes_RN\to Q$. Odwzorowania
$\alpha\beta$ i $\beta\alpha$ są identycznościami
na klasach tensorów prostych, które generują
odpowiednie moduły. Stąd $\beta$ jest izomorfizmem,
czyli $\ker(v\otimes1)=\operatorname{im}(u\otimes1)$.
\end{proof}
\begin{przyklad}[Dlaczego brak lewostronnej dokładności?]
Wstrzyknięcie $\Z\xrightarrow{\cdot2}\Z$ po zastosowaniu
funktora $-\otimes_\Z\Z/2\Z$ staje się mapą zerową
$\Z/2\Z\xrightarrow{0}\Z/2\Z$, która nie jest różnowartościowa.
Właśnie tę utratę informacji wykryje funktor $\operatorname{Tor}$.
\end{przyklad}

\begin{uwaga}[Ranga modułu nie jest liczbą generatorów]
Jeśli $R$ jest dziedziną z ciałem ułamków $K$, rangą $R$-modułu
$M$ nazywamy $\dim_K(K\otimes_R M)$; iloczyn tensorowy
został właśnie skonstruowany. Dla $R=\Z$
moduł $\Z/6\Z$ ma rangę $0$, choć potrzebuje generatora,
a $\Z^3\oplus\Z/2\Z$ ma rangę $3$.
Rangi $M$ nie należy utożsamiać z liczbą generatorów ani
z wymiarem dowolnego ilorazu $M/N$. Gdy $R$ jest ciałem,
ranga $M$ jest po prostu jego wymiarem jako przestrzeni wektorowej.
\end{uwaga}

\section{Projektywność, \texorpdfstring{$\operatorname{Tor}$ i $\operatorname{Ext}$}{Tor i Ext}}
W tym paragrafie zakładamy, że $R$ jest \emph{przemiennym}
pierścieniem z jedynką. Pozwala to stosować
iloczyn tensorowy dwóch lewych $R$-modułów.
Dla pierścieni nieprzemiennych trzeba osobno
rozróżnić moduły prawe i lewe.
\begin{definicja}
Moduł $P$ nad $R$ nazywamy \emph{projektywnym}, jeśli dla
każdego surjektywnego homomorfizmu $R$-modułów
$e:M\twoheadrightarrow N$ i każdego homomorfizmu
$R$-modułów $f:P\to N$ istnieje homomorfizm
$\widetilde f:P\to M$ z $e\widetilde f=f$.
Moduł $F$ nazywamy \emph{płaskim}, jeśli
$-\otimes_RF$ przekształca wszystkie krótkie ciągi dokładne
w krótkie ciągi dokładne.
\end{definicja}
\begin{stwierdzenie}
Niech $R$ będzie przemiennym pierścieniem z jedynką.
Każdy wolny $R$-moduł jest projektywny, a każdy
projektywny $R$-moduł jest płaski. Każdy $R$-moduł
jest ilorazem wolnego $R$-modułu.
\end{stwierdzenie}
\begin{proof}
Jeżeli $P$ jest wolny o bazie $(e_i)$, to dla każdego $f$
wybieramy preobrazy $f(e_i)$ przez surjekcję $e$ i przedłużamy
liniowo. To daje podniesienie. Dla dowolnego $M$ surjekcja
z wolnego modułu na generatorach symboli $[m]$ jest zadana
$[m]\mapsto m$. Moduł projektywny jest składnikiem prostym
wolnego: podnosząc $\id_P$ wzdłuż wolnej surjekcji $F\to P$
otrzymujemy $F\cong P\oplus\ker(F\to P)$.
Jeśli $F\cong R^{(I)}$ jest wolny, to dla dowolnego
$R$-modułu $M$ mamy naturalny izomorfizm
$M\otimes_RF\cong M^{(I)}$.
Dowodzimy tego, wysyłając $m\otimes e_i$ na
wektor o współrzędnej $m$ w pozycji $i$;
dwuliniowość daje mapę w jedną stronę,
a odwrotność jest sumą skończoną
$\sum_i m_i\otimes e_i$. W szczególności
dokładny ciąg po zastosowaniu $-\otimes_RF$
staje się sumą prostą kopii tego ciągu;
dokładność sprawdzamy współrzędnie.

Jeżeli $P$ jest projektywny, to
$F\cong P\oplus Q$ dla pewnego $Q$.
Dla dowolnego krótkiego ciągu dokładnego
ciąg otrzymany po zastosowaniu $-\otimes_RF$
rozkłada się na sumę prostą ciągów dla $P$ i $Q$.
Skoro całkowity ciąg jest dokładny, każdy z jego
dwóch składników również jest dokładny
(wynika to wprost z rozkładu jąder i obrazów).
Zatem $P$ jest płaski.
\end{proof}
\begin{definicja}[Wolna prezentacja i pierwsze funktory pochodne]
Dla grup abelowych $A,B$ wybieramy wolną prezentację
$0\to F_1\xrightarrow{d}F_0\xrightarrow{p}A\to0$,
gdzie $F_0,F_1$ są wolnymi grupami abelowymi.
Taka prezentacja istnieje, ponieważ każda podgrupa
wolnej grupy abelowej jest wolną grupą abelową.
Definiujemy
$\operatorname{Tor}_1^{\Z}(A,B)=\ker(d\otimes1_B)$ oraz
$\operatorname{Ext}^1_{\Z}(A,B)=\operatorname{coker}
\bigl(\operatorname{Hom}(F_0,B)\xrightarrow{d^*}
\operatorname{Hom}(F_1,B)\bigr)$,
gdzie $d^*(f)=f\circ d$ dla $f\in\operatorname{Hom}(F_0,B)$.
Niezależność od prezentacji wynika z porównania prezentacji
za pomocą podnoszeń homomorfizmów wolnych grup; dowód
zapiszemy w lemacie niżej.
\end{definicja}
\begin{lemat}[Niezależność od prezentacji]
Niech $A$ będzie grupą abelową, a
$0\to F_1\xrightarrow{d}F_0\xrightarrow{p}A\to0$
i $0\to F'_1\xrightarrow{d'}F'_0\xrightarrow{p'}A\to0$
dwiema wolnymi prezentacjami $A$.
Istnieją mapy łańcuchowe między ich dwustopniowymi
kompleksami indukujące $\id_A$ na homologiach stopnia
zerowego, a ich kompozycje są łańcuchowo homotopijne
z identycznościami. Dlatego $\operatorname{Tor}_1^\Z(A,B)$
i $\operatorname{Ext}^1_\Z(A,B)$ nie zależą od wyboru
wolnej prezentacji, dla każdej grupy abelowej $B$.
\end{lemat}
\begin{proof}
Niech $p:F_0\twoheadrightarrow A$, $p':F'_0\twoheadrightarrow A$.
Wybieramy dla każdego elementu bazy $F_0$ podniesienie przez
$p'$ jego obrazu przez $p$, uzyskując $f_0:F_0\to F'_0$
z $p'f_0=p$. Mapa $f_0d$ ląduje w $\ker p'=\operatorname{im}d'$, 
a $d':F'_1\to\operatorname{im}d'$ jest izomorfizmem;
można zatem położyć $f_1=(d')^{-1}f_0d$.
Analogicznie konstruujemy $g_0,g_1$ w przeciwną stronę.
Ponieważ $gf-\id$ indukuje zero na $A$, obraz
$g_0f_0-\id_{F_0}$ leży w $\operatorname{im}d$.
Definiujemy $h_0=d^{-1}(g_0f_0-\id)$ na $F_0$.
Wtedy $dh_0=g_0f_0-\id$ oraz, ponieważ $d$ jest
różnowartościowe, $h_0d=g_1f_1-\id$;
jest to homotopia łańcuchowa.
Podobnie dla $fg-\id$. Funktor $-\otimes_\Z B$ przenosi równość
$gf-\id=dh+hd$ na analogiczną równość dla
odwzorowań kompleksów z $B$;
zachowuje więc homotopie łańcuchowe.
Kontrawariantny funktor $\operatorname{Hom}(-,B)$
odwraca kierunki strzałek, a relacja homotopii
przechodzi na odpowiednią relację w kompleksach
kohomologicznych.
Stąd izomorfizmy odpowiednich grup homologii i kohomologii.
\end{proof}
\begin{przyklad}
Dla $A=\Z/m\Z$ prezentacja
$0\to\Z\xrightarrow{\cdot m}\Z\to A\to0$ daje
\[
\operatorname{Tor}_1^\Z(A,B)=\{b\in B:mb=0\},\qquad
\operatorname{Ext}^1_\Z(A,B)=B/mB.
\]
Jeśli $A$ jest wolna, oba funktory znikają, gdyż można
wybrać prezentację $F_1=0$.
Dla $A=\Z/12\Z$ i $B=\Z/18\Z$ otrzymujemy
$\operatorname{Tor}_1^\Z(A,B)
=\{[0],[3],[6],[9],[12],[15]\}\cong\Z/6\Z$.
Ponadto $12B=\{[0],[6],[12]\}$,
więc $\operatorname{Ext}^1_\Z(A,B)=B/12B
\cong\Z/6\Z$.
Otrzymaliśmy izomorficzne grupy na dwa różne sposoby,
lecz nie należy stąd wyciągać ogólnego wniosku
o naturalnej identyczności funktorów
$\operatorname{Tor}_1$ i $\operatorname{Ext}^1$.
\end{przyklad}
\begin{uwaga}[Połączenie z następnym rozdziałem]
Dalszy rozdział o algebrze homologicznej wprowadzi
kompleksy i krótkie ciągi dokładne kompleksów. Definicje
$\operatorname{Tor}$ i $\operatorname{Ext}$ wykorzystamy potem
w twierdzeniu Künnetha oraz twierdzeniu o współczynnikach
uniwersalnych. Opisane konstrukcje nie są zwykłymi iloczynami
kartezjańskimi: iloczyn tensorowy może nawet być
zerowy dla dwóch niezerowych czynników.
\end{uwaga}
\clearpage


